% Fontes (por nome de arquivo, para funcionar no Tectonic sem fontes do sistema)
\setmainfont{LibertinusSerif}[Extension=.otf, UprightFont=*-Regular, BoldFont=*-Bold, ItalicFont=*-Italic, BoldItalicFont=*-BoldItalic]
\setsansfont{LibertinusSans}[Extension=.otf, UprightFont=*-Regular, BoldFont=*-Bold, ItalicFont=*-Italic]
\setmathfont{LibertinusMath-Regular.otf}
\setmonofont{DejaVuSansMono}[Extension=.ttf, BoldFont=*-Bold, Scale=0.86]
\usepackage{newunicodechar}
\newunicodechar{☐}{\ensuremath{\square}}
\newfontfamily\simbolos{DejaVuSans}[Extension=.ttf]
\newunicodechar{✓}{{\simbolos ✓}}
% Página e títulos
\usepackage[a4paper, inner=2.6cm, outer=2.4cm, top=2.6cm, bottom=2.8cm]{geometry}
\usepackage{xcolor}
\definecolor{azul}{HTML}{1F5F99}
\usepackage{titlesec}
\titleformat{\chapter}[display]{\normalfont\sffamily\bfseries\color{azul}\raggedright}{}{0pt}{\huge}
\titlespacing*{\chapter}{0pt}{-10pt}{28pt}
\titleformat{\section}{\normalfont\sffamily\bfseries\Large\color{azul}}{}{0pt}{}
\titleformat{\subsection}{\normalfont\sffamily\bfseries\large}{}{0pt}{}
\titleformat{\subsubsection}{\normalfont\sffamily\bfseries\normalsize}{}{0pt}{}
\usepackage{fancyhdr}
\pagestyle{fancy}\fancyhf{}
\fancyhead[LE]{\small\sffamily\leftmark}\fancyhead[RO]{\small\sffamily\rightmark}
\fancyfoot[C]{\small\thepage}
\renewcommand{\headrulewidth}{0.3pt}
\renewcommand{\chaptermark}[1]{\markboth{#1}{#1}}
\renewcommand{\sectionmark}[1]{\markright{#1}}
% Legendas: o número já vem no texto da legenda ("Figura 1.2 — ...")
\usepackage{caption}
\captionsetup{labelformat=empty, font=small, labelfont=bf, justification=justified, singlelinecheck=false}
\captionsetup[table]{labelformat=empty}
% Figuras fixas no ponto do texto (como no Word): nenhuma figura sai do lugar nem se perde entre páginas
\usepackage{float}
\floatplacement{figure}{H}
% Caixas (citações em markdown) com filete lateral, que quebram entre páginas
\usepackage{framed}
\renewenvironment{quote}{\def\FrameCommand{{\color{azul!70}\vrule width 2.2pt}\hspace{8pt}}\MakeFramed{\advance\hsize-\width\FrameRestore}\small}{\endMakeFramed}
% Código: quebra de linha automática
\usepackage{fvextra}
\DefineVerbatimEnvironment{Highlighting}{Verbatim}{breaklines,breakanywhere,commandchars=\\\{\},fontsize=\small}
% Tabelas
\renewcommand{\arraystretch}{1.15}
\setlength{\emergencystretch}{3em}
\usepackage{microtype}
\raggedbottom
\usepackage{etoolbox}
\AtBeginEnvironment{longtable}{\small}
\usepackage{titling}
\pretitle{\begin{center}\vspace*{5cm}\sffamily\bfseries\color{azul}\fontsize{26}{32}\selectfont}
\posttitle{\par\end{center}\vspace{0.6cm}}
\predate{\begin{center}\vspace{6cm}\large}
\postdate{\par\end{center}}
\renewcommand{\maketitlehookd}{}
